% GENERATED FILE. Edit canonical lessons, then run npm run generate:books.
% canonical-source-hash: fnv1a64:4e51bf7a9f57d8d6
% canonical-lessons: KA-C354-samrajya, KA-C354-cakravarti, KA-C354-pattabhiseka, KA-C354-asthana, KA-C354-rajakumara, KA-R354-review-348-354

\chapter{An Empire and Its Emperor}
\label{ch:ka-an-empire-and-its-emperor}
\hlchaptermodality{354}

\begin{chapteropening}
\textbf{By the end of this chapter:} \emph{I can say an empire, an emperor, a coronation, a royal court and a prince.}

Complete the last lesson of chapter 354: I can say an empire, an emperor, a coronation, a royal court and a prince.
\end{chapteropening}

\section[\texorpdfstring{\kn{ಸಾಮ್ರಾಜ್ಯ} (sāmrājya)}{ಸಾಮ್ರಾಜ್ಯ (sāmrājya)}]{\kn{ಸಾಮ್ರಾಜ್ಯ} (sāmrājya) — an empire}

\label{lesson:KA-C354-samrajya}



\begin{quote}
Before the new one: say the Kannada for a heritage, then the Kannada for to pester.
\end{quote}

\subsection*{You'll want to know: \kn{ಸಾಮ್ರಾಜ್ಯ}}
\textbf{\kn{ಸಾಮ್ರಾಜ್ಯ}} — \emph{sāmrājya} — \textquotedblleft{}an empire\textquotedblright{}.

\begin{quote}
\kn{ಆ} \kn{ಸಾಮ್ರಾಜ್ಯ} \kn{ಸಮುದ್ರದಿಂದ} \kn{ಸಮುದ್ರದವರೆಗೆ} \kn{ಇತ್ತು}.
\end{quote}

\begin{quote}
\emph{ā sāmrājya samudradinda samudradavarege ittu.}
\end{quote}

\begin{quote}
That empire stretched from sea to sea.
\end{quote}

Plural: \kn{ಸಾಮ್ರಾಜ್ಯಗಳು}, \emph{sāmrājyagaḷu}.

\begin{cousinweb}
\textbf{\kn{ಸಾಮ್ರಾಜ್ಯ}} (\emph{sāmrājya}) is borrowed from Sanskrit \emph{sāmrājya}, \textquotedblleft{}complete sovereignty, empire\textquotedblright{} (Monier-Williams), from \emph{samrāj}, \textquotedblleft{}a supreme ruler\textquotedblright{}.
\end{cousinweb}

\begin{grammarlens}[title={What you've built}]
One more word for telling what happened.
\end{grammarlens}

\subsection*{Your turn}
\begin{itemize}
\raggedright
  \item \emph{Say it:} \emph{sāmrājya}
  \item \emph{Say it:} \emph{sāmrājya}, once more
  \item \emph{Say it:} \emph{ā sāmrājya samudradinda samudradavarege ittu.} — That empire stretched from sea to sea.
  \item \emph{Recall:} say the Kannada for a heritage, then the Kannada for to pester, then say \emph{sāmrājya} again
\end{itemize}

\begin{tcolorbox}[breakable,colback=teal!4,colframe=teal!35!black,title={Before you move on}]
What does \kn{ಸಾಮ್ರಾಜ್ಯ} mean? (\textquotedblleft{}An empire\textquotedblright{}.) Say it once more.
\end{tcolorbox}
\section[\texorpdfstring{\kn{ಚಕ್ರವರ್ತಿ} (cakravarti)}{ಚಕ್ರವರ್ತಿ (cakravarti)}]{\kn{ಚಕ್ರವರ್ತಿ} (cakravarti) — an emperor}

\label{lesson:KA-C354-cakravarti}



\begin{quote}
Before the new one: say the Kannada for to pester, then the Kannada for an empire.
\end{quote}

\subsection*{You'll want to know: \kn{ಚಕ್ರವರ್ತಿ}}
\textbf{\kn{ಚಕ್ರವರ್ತಿ}} — \emph{cakravarti} — \textquotedblleft{}an emperor\textquotedblright{}.

\begin{quote}
\kn{ಚಕ್ರವರ್ತಿ} \kn{ಹತ್ತು} \kn{ರಾಜ್ಯಗಳನ್ನು} \kn{ಆಳಿದನು}.
\end{quote}

\begin{quote}
\emph{cakravarti hattu rājyagaḷannu āḷidanu.}
\end{quote}

\begin{quote}
The emperor ruled ten kingdoms.
\end{quote}

Plural: \kn{ಚಕ್ರವರ್ತಿಗಳು}, \emph{cakravartigaḷu}.

A ruler over many kings, the head of a \kn{ಸಾಮ್ರಾಜ್ಯ}.

\begin{cousinweb}
\textbf{\kn{ಚಕ್ರವರ್ತಿ}} (\emph{cakravarti}) is borrowed from Sanskrit \emph{cakravartin}, \textquotedblleft{}a universal ruler\textquotedblright{} (CDIAL 4549), literally one whose chariot wheel, \emph{cakra}, rolls everywhere.
\end{cousinweb}

\begin{grammarlens}[title={What you've built}]
One more word for telling what happened.
\end{grammarlens}

\subsection*{Your turn}
\begin{itemize}
\raggedright
  \item \emph{Say it:} \emph{cakravarti}
  \item \emph{Say it:} \emph{cakravarti}, once more
  \item \emph{Say it:} \emph{cakravarti hattu rājyagaḷannu āḷidanu.} — The emperor ruled ten kingdoms.
  \item \emph{Recall:} say the Kannada for to pester, then the Kannada for an empire, then say \emph{cakravarti} again
\end{itemize}

\begin{tcolorbox}[breakable,colback=teal!4,colframe=teal!35!black,title={Before you move on}]
What does \kn{ಚಕ್ರವರ್ತಿ} mean? (\textquotedblleft{}An emperor\textquotedblright{}.) Say it once more.
\end{tcolorbox}
\section[\texorpdfstring{\kn{ಪಟ್ಟಾಭಿಷೇಕ} (paṭṭābhiṣēka)}{ಪಟ್ಟಾಭಿಷೇಕ (paṭṭābhiṣēka)}]{\kn{ಪಟ್ಟಾಭಿಷೇಕ} (paṭṭābhiṣēka) — a coronation}

\label{lesson:KA-C354-pattabhiseka}



\begin{quote}
Before the new one: say the Kannada for an empire, then the Kannada for an emperor.
\end{quote}

\subsection*{You'll want to know: \kn{ಪಟ್ಟಾಭಿಷೇಕ}}
\textbf{\kn{ಪಟ್ಟಾಭಿಷೇಕ}} — \emph{paṭṭābhiṣēka} — \textquotedblleft{}a coronation\textquotedblright{}.

\begin{quote}
\kn{ಹೊಸ} \kn{ರಾಜನ} \kn{ಪಟ್ಟಾಭಿಷೇಕ} \kn{ಅರಮನೆಯಲ್ಲಿ} \kn{ನಡೆಯಿತು}.
\end{quote}

\begin{quote}
\emph{hosa rājana paṭṭābhiṣēka aramaneyalli naḍeyitu.}
\end{quote}

\begin{quote}
The new king's coronation took place in the palace.
\end{quote}

\kn{ಪಟ್ಟ}, a crown band, and \kn{ಅಭಿಷೇಕ}, an anointing.

\begin{cousinweb}
\textbf{\kn{ಪಟ್ಟಾಭಿಷೇಕ}} (\emph{paṭṭābhiṣēka}) is borrowed from Sanskrit \emph{paṭṭābhiṣeka}, \textquotedblleft{}consecration of a tiara\textquotedblright{} (Monier-Williams), from \emph{paṭṭa} \textquotedblleft{}a band\textquotedblright{} and \emph{abhiṣeka} \textquotedblleft{}anointing\textquotedblright{}.
\end{cousinweb}

\begin{grammarlens}[title={What you've built}]
One more word for telling what happened.
\end{grammarlens}

\subsection*{Your turn}
\begin{itemize}
\raggedright
  \item \emph{Say it:} \emph{paṭṭābhiṣēka}
  \item \emph{Say it:} \emph{paṭṭābhiṣēka}, once more
  \item \emph{Say it:} \emph{hosa rājana paṭṭābhiṣēka aramaneyalli naḍeyitu.} — The new king's coronation took place in the palace.
  \item \emph{Recall:} say the Kannada for an empire, then the Kannada for an emperor, then say \emph{paṭṭābhiṣēka} again
\end{itemize}

\begin{tcolorbox}[breakable,colback=teal!4,colframe=teal!35!black,title={Before you move on}]
What does \kn{ಪಟ್ಟಾಭಿಷೇಕ} mean? (\textquotedblleft{}A coronation\textquotedblright{}.) Say it once more.
\end{tcolorbox}
\section[\texorpdfstring{\kn{ಆಸ್ಥಾನ} (āsthāna)}{ಆಸ್ಥಾನ (āsthāna)}]{\kn{ಆಸ್ಥಾನ} (āsthāna) — a royal court}

\label{lesson:KA-C354-asthana}



\begin{quote}
Before the new one: say the Kannada for an emperor, then the Kannada for a coronation.
\end{quote}

\subsection*{You'll want to know: \kn{ಆಸ್ಥಾನ}}
\textbf{\kn{ಆಸ್ಥಾನ}} — \emph{āsthāna} — \textquotedblleft{}a royal court\textquotedblright{}.

\begin{quote}
\kn{ರಾಜನ} \kn{ಆಸ್ಥಾನದಲ್ಲಿ} \kn{ಕವಿಗಳು} \kn{ಮತ್ತು} \kn{ಸಂಗೀತಗಾರರು} \kn{ಇದ್ದರು}.
\end{quote}

\begin{quote}
\emph{rājana āsthānadalli kavigaḷu mattu saṅgītagāraru iddaru.}
\end{quote}

\begin{quote}
There were poets and musicians at the king's court.
\end{quote}

The king's hall and the people around him: \kn{ಆಸ್ಥಾನ} \kn{ಸಂಗೀತಗಾರ}, a court musician.

\begin{cousinweb}
\textbf{\kn{ಆಸ್ಥಾನ}} (\emph{āsthāna}) is borrowed from Sanskrit \emph{āsthāna}, \textquotedblleft{}a place, an assembly\textquotedblright{} (CDIAL 1514).
\end{cousinweb}

\begin{grammarlens}[title={What you've built}]
One more word for telling what happened.
\end{grammarlens}

\subsection*{Your turn}
\begin{itemize}
\raggedright
  \item \emph{Say it:} \emph{āsthāna}
  \item \emph{Say it:} \emph{āsthāna}, once more
  \item \emph{Say it:} \emph{rājana āsthānadalli kavigaḷu mattu saṅgītagāraru iddaru.} — There were poets and musicians at the king's court.
  \item \emph{Recall:} say the Kannada for an emperor, then the Kannada for a coronation, then say \emph{āsthāna} again
\end{itemize}

\begin{tcolorbox}[breakable,colback=teal!4,colframe=teal!35!black,title={Before you move on}]
What does \kn{ಆಸ್ಥಾನ} mean? (\textquotedblleft{}A royal court\textquotedblright{}.) Say it once more.
\end{tcolorbox}
\section[\texorpdfstring{\kn{ರಾಜಕುಮಾರ} (rājakumāra)}{ರಾಜಕುಮಾರ (rājakumāra)}]{\kn{ರಾಜಕುಮಾರ} (rājakumāra) — a prince}

\label{lesson:KA-C354-rajakumara}



\begin{quote}
Before the new one: say the Kannada for a coronation, then the Kannada for a royal court.
\end{quote}

\subsection*{You'll want to know: \kn{ರಾಜಕುಮಾರ}}
\textbf{\kn{ರಾಜಕುಮಾರ}} — \emph{rājakumāra} — \textquotedblleft{}a prince\textquotedblright{}.

\begin{quote}
\kn{ರಾಜಕುಮಾರ} \kn{ಕುದುರೆ} \kn{ಏರಿ} \kn{ಕಾಡಿಗೆ} \kn{ಹೋದನು}.
\end{quote}

\begin{quote}
\emph{rājakumāra kudure ēri kāḍige hōdanu.}
\end{quote}

\begin{quote}
The prince rode a horse to the forest.
\end{quote}

Plural: \kn{ರಾಜಕುಮಾರರು}, \emph{rājakumāraru}.

A princess is a \kn{ರಾಜಕುಮಾರಿ}.

\begin{cousinweb}
\textbf{\kn{ರಾಜಕುಮಾರ}} (\emph{rājakumāra}) is borrowed from Sanskrit \emph{rājakumāra}, \textquotedblleft{}a king's son, a prince\textquotedblright{} (Monier-Williams), from \emph{rāja} \textquotedblleft{}king\textquotedblright{} and \emph{kumāra} \textquotedblleft{}a boy, a son\textquotedblright{}.
\end{cousinweb}

\begin{grammarlens}[title={What you've built}]
One more word for telling what happened.
\end{grammarlens}

\subsection*{Your turn}
\begin{itemize}
\raggedright
  \item \emph{Say it:} \emph{rājakumāra}
  \item \emph{Say it:} \emph{rājakumāra}, once more
  \item \emph{Say it:} \emph{rājakumāra kudure ēri kāḍige hōdanu.} — The prince rode a horse to the forest.
  \item \emph{Recall:} say the Kannada for a coronation, then the Kannada for a royal court, then say \emph{rājakumāra} again
\end{itemize}

\begin{tcolorbox}[breakable,colback=teal!4,colframe=teal!35!black,title={Before you move on}]
What does \kn{ರಾಜಕುಮಾರ} mean? (\textquotedblleft{}A prince\textquotedblright{}.) Say it once more.
\end{tcolorbox}
\section[(dialogue)]{Looking back: from The Pirate Treasure to An Empire and Its Emperor}

\label{lesson:KA-R354-review-348-354}



\begin{quote}
Say the Kannada for a treasure and a pirate.
\end{quote}

\subsection*{Your turn}
Say these aloud:

\begin{itemize}
\raggedright
  \item a treasure, a pirate, an adventure, a secret, a sorcerer
  \item a sword, an arrow, armour, a charioteer, to fall in battle
  \item to honour, to extol, an honorific title, fame, to inaugurate
  \item a mission, a volunteer, to crash against, to engulf, to be washed away
  \item a revolt, a revolution, a slave, to set free, to capture
  \item a will, an heir, a legacy, a heritage, to pester
  \item an empire, an emperor, a coronation, a royal court, a prince
\end{itemize}

\begin{tcolorbox}[breakable,colback=teal!4,colframe=teal!35!black,title={Before you move on}]
A treasure? (\textbf{\kn{ನಿಧಿ}}, \emph{nidhi}.) A pirate? (\textbf{\kn{ಕಡಲ್ಗಳ್ಳ}}, \emph{kaḍalgaḷḷa}.) An adventure? (\textbf{\kn{ಸಾಹಸ}}, \emph{sāhasa}.) A secret? (\textbf{\kn{ರಹಸ್ಯ}}, \emph{rahasya}.) A sorcerer? (\textbf{\kn{ಮಾಂತ್ರಿಕ}}, \emph{māntrika}.) A sword? (\textbf{\kn{ಖಡ್ಗ}}, \emph{khaḍga}.) An arrow? (\textbf{\kn{ಬಾಣ}}, \emph{bāṇa}.) Armour? (\textbf{\kn{ಕವಚ}}, \emph{kavaca}.) A charioteer? (\textbf{\kn{ಸಾರಥಿ}}, \emph{sārathi}.) To fall in battle? (\textbf{\kn{ಮಡಿ}}, \emph{maḍi}.) To honour? (\textbf{\kn{ಸನ್ಮಾನಿಸು}}, \emph{sanmānisu}.) To extol? (\textbf{\kn{ಕೊಂಡಾಡು}}, \emph{koṇḍāḍu}.) An honorific title? (\textbf{\kn{ಬಿರುದು}}, \emph{birudu}.) Fame? (\textbf{\kn{ಕೀರ್ತಿ}}, \emph{kīrti}.) To inaugurate? (\textbf{\kn{ಉದ್ಘಾಟಿಸು}}, \emph{udghāṭisu}.) A mission? (\textbf{\kn{ಕಾರ್ಯಾಚರಣೆ}}, \emph{kāryācaraṇe}.) A volunteer? (\textbf{\kn{ಸ್ವಯಂಸೇವಕ}}, \emph{svayaṃsēvaka}.) To crash against? (\textbf{\kn{ಅಪ್ಪಳಿಸು}}, \emph{appaḷisu}.) To engulf? (\textbf{\kn{ಆವರಿಸು}}, \emph{āvarisu}.) To be washed away? (\textbf{\kn{ಕೊಚ್ಚಿಹೋಗು}}, \emph{koccihōgu}.) A revolt? (\textbf{\kn{ದಂಗೆ}}, \emph{daṅge}.) A revolution? (\textbf{\kn{ಕ್ರಾಂತಿ}}, \emph{krānti}.) A slave? (\textbf{\kn{ಗುಲಾಮ}}, \emph{gulāma}.) To set free? (\textbf{\kn{ಬಿಡುಗಡೆ} \kn{ಮಾಡು}}, \emph{biḍugaḍe māḍu}.) To capture? (\textbf{\kn{ವಶಪಡಿಸಿಕೊ}}, \emph{vaśapaḍisiko}.) A will? (\textbf{\kn{ಉಯಿಲು}}, \emph{uyilu}.) An heir? (\textbf{\kn{ವಾರಸುದಾರ}}, \emph{vārasudāra}.) A legacy? (\textbf{\kn{ಬಳುವಳಿ}}, \emph{baḷuvaḷi}.) A heritage? (\textbf{\kn{ಪರಂಪರೆ}}, \emph{parampare}.) To pester? (\textbf{\kn{ಪೀಡಿಸು}}, \emph{pīḍisu}.) An empire? (\textbf{\kn{ಸಾಮ್ರಾಜ್ಯ}}, \emph{sāmrājya}.) An emperor? (\textbf{\kn{ಚಕ್ರವರ್ತಿ}}, \emph{cakravarti}.) A coronation? (\textbf{\kn{ಪಟ್ಟಾಭಿಷೇಕ}}, \emph{paṭṭābhiṣēka}.) A royal court? (\textbf{\kn{ಆಸ್ಥಾನ}}, \emph{āsthāna}.) A prince? (\textbf{\kn{ರಾಜಕುಮಾರ}}, \emph{rājakumāra}.)
\end{tcolorbox}
